% Lean: EconLean/Macro/OLG/{LogHousehold,Diamond,Samuelson}.lean
\chapter{Overlapping Generations}
\label{chap:olg}

\section{The two-period log household}
\label{sec:log_household}

Lifetime utility is $U(c_1, c_2) = \log c_1 + \beta \log c_2$ with $\beta > 0$; the budget is
$c_1 + c_2/R \le W$ with $R, W > 0$.

\begin{definition}[Log lifetime utility]
\label{def:logU}
\lean{Econ.OLG.logU}
\leanok
$\mathrm{logU}_\beta(c_1, c_2) = \log c_1 + \beta \log c_2$.
\end{definition}

\begin{theorem}[Optimal consumption]
\label{thm:logU_opt}
\lean{Econ.OLG.logU_le_opt, Econ.OLG.logU_lt_opt}
\leanok
\uses{def:logU}
Let $c_1^* = W/(1+\beta)$, $c_2^* = \beta R W/(1 + \beta)$. For every affordable $c_1, c_2 > 0$,
$U(c_1, c_2) \le U(c_1^*, c_2^*)$, with strict inequality unless $(c_1, c_2) = (c_1^*, c_2^*)$.
\end{theorem}
\begin{proof}
\leanok
No calculus: apply $\log x \le x - 1$ to $x = c_1/c_1^*$ and $x = c_2/c_2^*$, weighting the second by $\beta$:
\[
  U(c_1, c_2) - U(c_1^*, c_2^*) \le \frac{c_1}{c_1^*} + \beta\frac{c_2}{c_2^*} - (1 + \beta).
\]
By algebra $\frac{c_1}{c_1^*} + \beta \frac{c_2}{c_2^*} = \frac{(1+\beta)(c_1 + c_2/R)}{W} \le 1 + \beta$.
For strictness, if $c_1 \neq c_1^*$ (or $c_2 \neq c_2^*$) the corresponding inequality $\log x < x - 1$ is strict.
\end{proof}

\begin{definition}[Optimal saving]
\label{def:isOptimalSaving}
\lean{Econ.OLG.IsOptimalSaving}
\leanok
\uses{def:logU}
With endowments $(\omega_1, \omega_2)$ and return $R$, saving $a$ is optimal if $\omega_1 - a > 0$,
$\omega_2 + R a > 0$, and no $a'$ with positive consumption gives higher $U(\omega_1 - a', \omega_2 + R a')$.
\end{definition}

\begin{theorem}[Optimal saving with log utility]
\label{thm:saving_isOptimal_iff}
\lean{Econ.OLG.saving_isOptimal_iff}
\leanok
\uses{def:isOptimalSaving, thm:logU_opt}
For $\beta, R, \omega_1 > 0$ and $\omega_2 \ge 0$: $a$ is optimal iff $a = (\beta\omega_1 - \omega_2/R)/(1 + \beta)$.
\end{theorem}
\begin{proof}
\leanok
Let $W = \omega_1 + \omega_2/R > 0$ and $a^*$ the stated value. Algebra gives $\omega_1 - a^* = c_1^*$ and
$\omega_2 + R a^* = c_2^*$, and every $a'$ satisfies $(\omega_1 - a') + (\omega_2 + R a')/R = W$. ($\Leftarrow$)
Theorem~\ref{thm:logU_opt}. ($\Rightarrow$) If $a \neq a^*$ then $\omega_1 - a \neq c_1^*$, so $a$ is strictly
worse than $a^*$ by Theorem~\ref{thm:logU_opt}, contradicting optimality.
\end{proof}

\section{The Diamond model}
\label{sec:diamond}

Sources: Diamond (1965); Acemoglu \S9.2--9.3. Log utility, $y = k^\alpha$, full depreciation, population
growth $n$. Wages $w(k) = (1-\alpha)k^\alpha$ and returns $R(k) = \alpha k^{\alpha-1}$.

\begin{definition}[Diamond economy]
\label{def:diamond}
\lean{Econ.OLG.Diamond, Econ.OLG.Diamond.wage, Econ.OLG.Diamond.rate, Econ.OLG.Diamond.coef, Econ.OLG.Diamond.next, Econ.OLG.Diamond.coef_pos, Econ.OLG.Diamond.wage_pos, Econ.OLG.Diamond.rate_pos, Econ.OLG.Diamond.next_pos}
\leanok
Parameters $0 < \alpha < 1$, $\beta > 0$, $n > -1$. Let $A = \frac{\beta(1-\alpha)}{(1+\beta)(1+n)} > 0$ and
$\mathrm{next}(k) = A k^\alpha$. Wage, return and $\mathrm{next}$ are positive at $k > 0$.
\end{definition}

\begin{definition}[Competitive equilibrium]
\label{def:diamond_eq}
\lean{Econ.OLG.Diamond.IsEquilibrium}
\leanok
\uses{def:diamond, def:isOptimalSaving}
A path $(k_t)$ is an equilibrium if $k_t > 0$ and, for every $t$, saving $(1+n)k_{t+1}$ is optimal for a young
agent with endowments $(w(k_t), 0)$ facing return $R(k_{t+1})$.
\end{definition}

\begin{theorem}[Equilibrium iff law of motion]
\label{thm:diamond_isEquilibrium_iff}
\lean{Econ.OLG.Diamond.isEquilibrium_iff}
\leanok
\uses{def:diamond_eq, thm:saving_isOptimal_iff}
For $k_0 > 0$, $(k_t)$ is an equilibrium iff $k_{t+1} = A k_t^\alpha$ for all $t$.
\end{theorem}
\begin{proof}
\leanok
At positive $k_t, k_{t+1}$, by Theorem~\ref{thm:saving_isOptimal_iff} with $\omega_2 = 0$, optimality of
$(1+n)k_{t+1}$ means $(1+n)k_{t+1} = \frac{\beta}{1+\beta}(1-\alpha)k_t^\alpha$, i.e.\ $k_{t+1} = A k_t^\alpha$
(the interest rate drops out). ($\Rightarrow$) apply this at each $t$. ($\Leftarrow$) positivity of all $k_t$
follows by induction since $\mathrm{next}$ preserves positivity; then apply the equivalence.
\end{proof}

\begin{definition}[Diamond as a Solow economy]
\label{def:diamond_toSolow}
\lean{Econ.OLG.Diamond.toSolow, Econ.OLG.Diamond.toSolow_next}
\leanok
\uses{def:diamond, def:solow_economy, thm:solow_cobbDouglas}
Take $f(k) = k^\alpha$, $s = \beta(1-\alpha)/(1+\beta) \in (0, 1]$, $\delta = 1$ and the same $n$. Its law of
motion is $\frac{s k^\alpha}{1+n} = A k^\alpha$, the Diamond law of motion.
\end{definition}

\begin{theorem}[Existence and uniqueness of equilibrium]
\label{thm:diamond_existsUnique}
\lean{Econ.OLG.Diamond.existsUnique_equilibrium}
\leanok
\uses{thm:diamond_isEquilibrium_iff}
From any $k_0 > 0$ there is exactly one equilibrium path.
\end{theorem}
\begin{proof}
\leanok
The iterates $\mathrm{next}^t(k_0)$ satisfy the law of motion, so form an equilibrium
(Theorem~\ref{thm:diamond_isEquilibrium_iff}); any equilibrium from $k_0$ satisfies the same recursion, so
equals it by induction.
\end{proof}

\begin{theorem}[Steady state and global stability; Acemoglu Prop.~9.3]
\label{thm:diamond_tendsto}
\lean{Econ.OLG.Diamond.kstar, Econ.OLG.Diamond.kstar_pos, Econ.OLG.Diamond.isSteadyState_kstar, Econ.OLG.Diamond.next_kstar, Econ.OLG.Diamond.tendsto_equilibrium}
\leanok
\uses{def:diamond_toSolow, thm:diamond_isEquilibrium_iff, thm:solow_tendsto, thm:solow_cobbDouglas}
$k^* = A^{1/(1-\alpha)}$ is the steady state, and every equilibrium path from $k_0 > 0$ converges to it.
\end{theorem}
\begin{proof}
\leanok
In the Solow representation $s/(n + \delta) = A$, so the Cobb--Douglas closed form
(Theorem~\ref{thm:solow_cobbDouglas}) gives $k^*$. An equilibrium path coincides with the Solow path (induction
with Definition~\ref{def:diamond_toSolow}), which converges by Theorem~\ref{thm:solow_tendsto}.
\end{proof}

\begin{lemma}[Steady-state interest rate]
\label{lem:diamond_rate_steadyState}
\lean{Econ.OLG.Diamond.kstar_rpow, Econ.OLG.Diamond.rate_steadyState}
\leanok
\uses{thm:diamond_tendsto}
$(k^*)^{\alpha - 1} = A^{-1}$, hence $R(k^*) = \alpha/A = \frac{\alpha(1+\beta)(1+n)}{\beta(1-\alpha)}$.
\end{lemma}
\begin{proof}
\leanok
$(A^{1/(1-\alpha)})^{\alpha-1} = A^{-1}$; substitute into $R(k) = \alpha k^{\alpha-1}$.
\end{proof}

\begin{lemma}[Golden-rule capital]
\label{lem:diamond_kgold}
\lean{Econ.OLG.Diamond.kgold, Econ.OLG.Diamond.kgold_pos, Econ.OLG.Diamond.hasDerivAt_kgold, Econ.OLG.Diamond.rate_kgold, Econ.OLG.Diamond.rate_strictAntiOn}
\leanok
\uses{thm:cobbDouglas_goldenRule}
$k_g = (\alpha/(1+n))^{1/(1-\alpha)} > 0$ satisfies $f'(k_g) = R(k_g) = 1 + n$, and $R$ is strictly decreasing
on $(0, \infty)$.
\end{lemma}
\begin{proof}
\leanok
Theorem~\ref{thm:cobbDouglas_goldenRule} with $\delta = 1$; uniqueness of derivatives identifies $f'(k_g)$
with $R(k_g)$. $k^{\alpha-1}$ is strictly decreasing since $\alpha - 1 < 0$.
\end{proof}

\begin{theorem}[Dynamic inefficiency criterion; Acemoglu Prop.~9.4]
\label{thm:diamond_dynamicallyInefficient_iff}
\lean{Econ.OLG.Diamond.dynamicallyInefficient_iff}
\leanok
\uses{lem:diamond_rate_steadyState, lem:diamond_kgold}
The following are equivalent: $R(k^*) < 1 + n$; $\alpha(1+\beta) < \beta(1-\alpha)$; $k_g < k^*$.
\end{theorem}
\begin{proof}
\leanok
First equivalence: by Lemma~\ref{lem:diamond_rate_steadyState}, multiply through by $\beta(1-\alpha) > 0$ and cancel
$1 + n > 0$. Second: $R(k_g) = 1 + n$ and $R$ is strictly decreasing, so $R(k^*) < R(k_g) \iff k_g < k^*$.
\end{proof}

\begin{definition}[Feasible paths]
\label{def:diamond_feasible}
\lean{Econ.OLG.Diamond.IsFeasible, Econ.OLG.Diamond.isFeasible_steadyState}
\leanok
\uses{def:diamond}
$(k_t, C_t)$ is feasible if $k_t \ge 0$ and $C_t + (1+n)k_{t+1} \le k_t^\alpha$. The competitive steady state
with $C^* = (k^*)^\alpha - (1+n)k^*$ is feasible.
\end{definition}

\begin{theorem}[Overaccumulation is inefficient]
\label{thm:diamond_exists_dominating_path}
\lean{Econ.OLG.Diamond.exists_dominating_path}
\leanok
\uses{thm:diamond_dynamicallyInefficient_iff, def:diamond_feasible, thm:goldenRule}
If $R(k^*) < 1 + n$, there is a feasible path from $k^*$ with $C_t > C^*$ for every $t$.
\end{theorem}
\begin{proof}
\leanok
By Theorem~\ref{thm:diamond_dynamicallyInefficient_iff}, $k_g < k^*$. Set $k_0 = k^*$, $k_t = k_g$ for $t \ge 1$;
$C_0 = (k^*)^\alpha - (1+n)k_g$ and $C_t = k_g^\alpha - (1+n)k_g$ for $t \ge 1$. Feasibility holds with equality.
$C_0 > C^*$ because $k_g < k^*$; and $C_t > C^*$ for $t \ge 1$ by the golden rule
(Theorem~\ref{thm:goldenRule}) since $k^* \neq k_g$.
\end{proof}

\begin{lemma}[Steady-state consumption rises below the golden rule]
\label{lem:diamond_consumption_strictMono}
\lean{Econ.OLG.Diamond.consumption_strictMonoOn}
\leanok
\uses{lem:diamond_kgold, thm:goldenRule}
If $0 \le a < b \le k_g$, then $a^\alpha - (1+n)a < b^\alpha - (1+n)b$.
\end{lemma}
\begin{proof}
\leanok
If $b = k_g$, Theorem~\ref{thm:goldenRule}. Otherwise write $b = \lambda a + (1-\lambda)k_g$ with
$\lambda = (k_g - b)/(k_g - a) \in (0, 1)$. Strict concavity of $k^\alpha$ gives
$b^\alpha > \lambda a^\alpha + (1-\lambda)k_g^\alpha$, and $(1+n)b$ is the same combination of $(1+n)a$ and
$(1+n)k_g$. So $g(b) > \lambda g(a) + (1-\lambda) g(k_g) \ge g(a)$, using $g(k_g) \ge g(a)$
(Theorem~\ref{thm:goldenRule}), where $g(k) = k^\alpha - (1+n)k$.
\end{proof}

\begin{theorem}[Dynamic efficiency when $R^* \ge 1 + n$]
\label{thm:diamond_no_dominating_path}
\lean{Econ.OLG.Diamond.no_dominating_path}
\leanok
\uses{lem:lt_tangent, lem:diamond_consumption_strictMono, thm:diamond_dynamicallyInefficient_iff, def:diamond_feasible}
If $1 + n \le R(k^*)$, $k_0 = k^*$, $(k, C)$ is feasible and $C_t \ge C^*$ for all $t$, then $C_t = C^*$ for all $t$.
\end{theorem}
\begin{proof}
\leanok
Write $c = 1 + n$, $R^* = R(k^*)$.
\emph{Tangent bound.} By concavity, $k^\alpha \le (k^*)^\alpha + R^*(k - k^*)$ for $k \ge 0$
(Lemma~\ref{lem:lt_tangent}). With feasibility,
$c(k_{t+1} - k^*) \le R^*(k_t - k^*) - (C_t - C^*)$.
\emph{Capital falls weakly and stays below $k^*$.} By induction, $k_t \le k^*$ and $k_{t+1} \le k_t$: if
$k_t \le k^*$ then $R^*(k_t - k^*) \le c(k_t - k^*)$ (as $R^* \ge c$), and $C_t \ge C^*$ gives
$c\, k_{t+1} \le c\, k_t$.
\emph{A strict gain pushes capital strictly below $k^*$.} Suppose $C_s > C^*$. Since $R^*(k_s - k^*) \le 0$,
we get $k_{s+1} < k^*$, and $k_{s+1+N} \le k_{s+1}$ for all $N$.
\emph{A fixed per-period loss.} $k^* \le k_g$ (else $R^* < c$ by Theorem~\ref{thm:diamond_dynamicallyInefficient_iff}).
Let $\mathrm{gap} = g(k^*) - g(k_{s+1}) > 0$ (Lemma~\ref{lem:diamond_consumption_strictMono}). For
$t = s + 1 + N$, $g(k_t) \le g(k_{s+1})$ since $k_t \le k_{s+1} \le k_g$, and feasibility with $C_t \ge C^*$
gives $c(k_t - k_{t+1}) \ge C_t - g(k_t) + \ldots \ge \mathrm{gap}$. Inductively
$N \cdot \mathrm{gap} \le c(k_{s+1} - k_{s+1+N})$.
\emph{Contradiction.} Choose $N > c k^*/\mathrm{gap}$; then $k_{s+1+N} < k_{s+1} - k^* < 0$, contradicting $k \ge 0$.
\end{proof}

\begin{theorem}[Complete characterization; Acemoglu Prop.~9.4]
\label{thm:diamond_iff_improvement}
\lean{Econ.OLG.Diamond.dynamicallyInefficient_iff_exists_improvement}
\leanok
\uses{thm:diamond_exists_dominating_path, thm:diamond_no_dominating_path}
$R(k^*) < 1 + n$ iff some feasible path from $k^*$ has $C_t \ge C^*$ for all $t$ and $C_t > C^*$ for some $t$.
\end{theorem}
\begin{proof}
\leanok
($\Rightarrow$) Theorem~\ref{thm:diamond_exists_dominating_path}. ($\Leftarrow$) otherwise
Theorem~\ref{thm:diamond_no_dominating_path} forces $C_t = C^*$, contradicting the strict gain. The Lean file
also checks that both cases occur ($\alpha = 1/4, \beta = 1, n = 0$ is inefficient; $\alpha = 1/2$ is efficient).
\end{proof}

\section{Samuelson's exchange economy and fiat money}
\label{sec:samuelson}

Sources: Samuelson (1958); Shell (1971); Ljungqvist--Sargent ch.~9. One agent is born each period, lives two
periods with endowments $(\omega_1, \omega_2)$ and log utility; there is an initial old agent at $t = 0$.

\begin{definition}[Economy, allocations, feasibility, dominance]
\label{def:samuelson}
\lean{Econ.OLG.Samuelson, Econ.OLG.Samuelson.Allocation, Econ.OLG.Samuelson.IsFeasible, Econ.OLG.Samuelson.utility, Econ.OLG.Samuelson.ParetoDominates, Econ.OLG.Samuelson.StrictlyDominates, Econ.OLG.Samuelson.StrictlyDominates.paretoDominates}
\leanok
\uses{def:logU}
$\omega_1, \omega_2, \beta > 0$. An allocation gives consumption $\mathrm{young}_t$ to generation $t$ when young
and $\mathrm{old}_t$ to the old at date $t$ ($\mathrm{old}_0$ is the initial old). It is feasible if consumption is
positive and $\mathrm{young}_t + \mathrm{old}_t \le \omega_1 + \omega_2$. Generation $t$'s utility is
$U(\mathrm{young}_t, \mathrm{old}_{t+1})$. $x$ \emph{Pareto-dominates} $y$ if nobody is worse off and someone is
strictly better off; $x$ \emph{strictly dominates} $y$ if the initial old and every generation are strictly
better off (which implies Pareto dominance).
\end{definition}

\begin{definition}[Autarky]
\label{def:autarky}
\lean{Econ.OLG.Samuelson.autarky, Econ.OLG.Samuelson.autarky_feasible, Econ.OLG.Samuelson.autarkyRate, Econ.OLG.Samuelson.autarkyRate_lt_one_iff}
\leanok
\uses{def:samuelson}
Autarky gives everyone their endowment; it is feasible. The autarky rate is $R_{\mathrm{aut}} = \omega_2/(\beta\omega_1)$,
and $R_{\mathrm{aut}} < 1 \iff \omega_2 < \beta\omega_1$.
\end{definition}

\begin{theorem}[Equilibrium interest rates]
\label{thm:samuelson_isEquilibrium_iff}
\lean{Econ.OLG.Samuelson.IsEquilibrium, Econ.OLG.Samuelson.isEquilibrium_iff}
\leanok
\uses{def:autarky, thm:saving_isOptimal_iff}
A rate path $R$ supports zero saving as optimal at every date iff $R_t = R_{\mathrm{aut}}$ for all $t$.
\end{theorem}
\begin{proof}
\leanok
By Theorem~\ref{thm:saving_isOptimal_iff}, $0$ is optimal at $R_t$ iff $0 = (\beta\omega_1 - \omega_2/R_t)/(1+\beta)$,
i.e.\ $R_t = \omega_2/(\beta\omega_1)$.
\end{proof}

\begin{theorem}[The competitive equilibrium is autarky]
\label{thm:samuelson_ce_iff}
\lean{Econ.OLG.Samuelson.IsCompetitiveEquilibrium, Econ.OLG.Samuelson.isCompetitiveEquilibrium_iff}
\leanok
\uses{thm:samuelson_isEquilibrium_iff}
A nonmonetary competitive equilibrium $(R, x)$ (optimal savings $a_t$, market clearing $a_t = 0$, consumption
from the budget constraints, $\mathrm{old}_0 = \omega_2$) holds iff $R_t = R_{\mathrm{aut}}$ for all $t$ and $x$ is autarky.
\end{theorem}
\begin{proof}
\leanok
($\Rightarrow$) Market clearing gives $a_t = 0$, so Theorem~\ref{thm:samuelson_isEquilibrium_iff} pins down $R$, and the
budget constraints give $\mathrm{young}_t = \omega_1$, $\mathrm{old}_{t+1} = \omega_2$, $\mathrm{old}_0 = \omega_2$.
($\Leftarrow$) take $a = 0$.
\end{proof}

\begin{theorem}[Golden rule dominates autarky]
\label{thm:golden_dominates_autarky}
\lean{Econ.OLG.Samuelson.golden, Econ.OLG.Samuelson.golden_feasible, Econ.OLG.Samuelson.golden_dominates_autarky}
\leanok
\uses{def:autarky, thm:logU_opt}
The golden allocation splits $\omega_1 + \omega_2$ as $\frac{1}{1+\beta}$ to the young and $\frac{\beta}{1+\beta}$ to the old;
it is feasible. If $\omega_2 < \beta\omega_1$, it strictly dominates autarky.
\end{theorem}
\begin{proof}
\leanok
Initial old: $\omega_2 < \beta(\omega_1 + \omega_2)/(1+\beta)$ rearranges to $\omega_2 < \beta\omega_1$. Each generation: the
golden bundle is the unique optimum of Theorem~\ref{thm:logU_opt} with $R = 1$ and $W = \omega_1 + \omega_2$; the autarky
bundle is affordable and differs from it (since $\omega_1 \neq (\omega_1 + \omega_2)/(1+\beta)$ when $\omega_2 < \beta\omega_1$),
so it is strictly worse.
\end{proof}

\begin{theorem}[The First Welfare Theorem fails]
\label{thm:firstWelfareTheorem_fails}
\lean{Econ.OLG.Samuelson.firstWelfareTheorem_fails, Econ.OLG.Samuelson.competitiveEquilibrium_dominated}
\leanok
\uses{thm:samuelson_ce_iff, thm:golden_dominates_autarky}
If $\omega_2 < \beta\omega_1$, a competitive equilibrium exists, and every competitive equilibrium allocation is strictly
dominated by a feasible allocation.
\end{theorem}
\begin{proof}
\leanok
By Theorem~\ref{thm:samuelson_ce_iff} the equilibrium allocation is autarky (and autarky at $R_{\mathrm{aut}}$ is one);
the golden allocation dominates it (Theorem~\ref{thm:golden_dominates_autarky}).
\end{proof}

\begin{theorem}[Stationary monetary equilibrium]
\label{thm:samuelson_monetary}
\lean{Econ.OLG.Samuelson.IsMonetaryEquilibrium, Econ.OLG.Samuelson.monetaryAllocation, Econ.OLG.Samuelson.isMonetaryEquilibrium_iff}
\leanok
\uses{thm:saving_isOptimal_iff}
With money stock $M > 0$ and constant price level $p > 0$, buying real balances $M/p$ (gross return $1$) is optimal iff
$\omega_2 < \beta\omega_1$ and $p = M(1+\beta)/(\beta\omega_1 - \omega_2)$.
\end{theorem}
\begin{proof}
\leanok
By Theorem~\ref{thm:saving_isOptimal_iff} with $R = 1$, optimality means $M/p = (\beta\omega_1 - \omega_2)/(1+\beta)$. Since
$M/p > 0$, this forces $\beta\omega_1 - \omega_2 > 0$, and solving gives $p$. Conversely that $p$ is positive and satisfies
the equation.
\end{proof}

\begin{theorem}[Money is a Pareto improvement]
\label{thm:monetary_dominates}
\lean{Econ.OLG.Samuelson.monetaryEquilibrium_eq_golden_and_dominates}
\leanok
\uses{thm:samuelson_monetary, thm:golden_dominates_autarky}
In a stationary monetary equilibrium the allocation is the golden allocation, which is feasible and strictly
dominates autarky.
\end{theorem}
\begin{proof}
\leanok
Substituting $M/p = (\beta\omega_1 - \omega_2)/(1+\beta)$, $\omega_1 - M/p = (\omega_1 + \omega_2)/(1+\beta)$ and
$\omega_2 + M/p = \beta(\omega_1 + \omega_2)/(1+\beta)$. Then apply Theorem~\ref{thm:golden_dominates_autarky}.
\end{proof}
